\section{The exact scalar envelope at a strict cutoff}\label{sec:scalar-envelope}
This section sharpens the elementary sufficient estimate above. It is
useful for explicit bounds, but scalar optimality alone does not settle
the spectral question. Section~\ref{sec:universal-tail} supplies that step.
Throughout these two sections fix $0<\rho<1$, $s=(1+\rho)/(1-\rho)$,
$k_s=(s-1)/s<K<1$, and put
\[
 a=1-K,\quad p_*=(s-1)a/K,\quad C_s=\frac{s^s}{(s-1)^{s-1}},
 \quad g=C_saK^{s-1},\quad h(p)=\frac{(a+Kp)^s}{p^{s-1}}.
\]
Then $0<p_*<1$ and $0<g<1$. For $\tau\ge g$ let $p_\tau$ be the
unique solution of $h(p_\tau)=\tau$ in $(0,p_*]$, and set
$U(\tau)=\tau+K^s(1-p_\tau)$.

\begin{theorem}\label{thm:scalar-envelope}
Among nonnegative random variables on arbitrary probability spaces with
$\mathbb Er=1$ and $\mathbb E[r^s\mathbf1_{r>K}]\le\tau$, the class is
empty for $\tau<g$. For $\tau\ge g$ the maximum of $\mathbb Er^s$ is
$U(\tau)$, attained by the law with value $K$ of probability $1-p_\tau$
and value $K+a/p_\tau$ of probability $p_\tau$.
At $\tau=g$ this is the unique admissible law; a uniform $L$-cube realizes
it precisely when $2^Lp_*$ is integral. For $\tau>g$, the supremum over
all finite uniform Boolean cubes is $U(\tau)$, approached by positive
mean-one densities with floor $K$ and strict tail at most $\tau$.
\end{theorem}
\begin{proof}
Logarithmic differentiation gives
\[
 \frac{h'(p)}{h(p)}=\frac{Kp-(s-1)a}{p(a+Kp)}.
\]
Thus $h$ decreases from infinity to $h(p_*)=g$ on $(0,p_*]$, then
increases to $h(1)=1$. For an admissible $r$, write
$p=\mathbb P(r>K)>0$ and $m=\mathbb E[r\mathbf1_{r>K}]$.
The complementary mean is at most $K(1-p)$, so $m\ge a+Kp$.
Conditional Jensen yields
\[
 \tau\ge\mathbb E[r^s\mathbf1_{r>K}]
 \ge m^s/p^{s-1}\ge h(p)\ge g.
\]
It also forces $p\ge p_\tau$. Consequently
\[
 \mathbb Er^s\le\tau+K^s(1-p)\le U(\tau).
\]
The stated two-point law has mean one, strict tail $h(p_\tau)=\tau$
and total moment $U(\tau)$, proving attainment. If $\tau=g$, every
inequality above is equality: $p=p_*$, the lower value is almost surely
$K$, and strict convexity forces a constant upper value. This proves
uniqueness and the finite-cube integrality condition.
For $\tau>g$, choose
$p_L=(\lfloor2^Lp_\tau\rfloor+1)/2^L$. Eventually
$p_\tau<p_L<p_*$ and $p_L\to p_\tau$. Assign value $K+a/p_L$ to
exactly $2^Lp_L$ vertices and $K$ to the others. Its mean is one and its
strict tail $h(p_L)<\tau$, whereas its total moment tends to $U(\tau)$.
Exact attainment on a specified cube requires the high probability
$p_\tau$ to be realizable there; no irrational probability is assigned
to a finite uniform event.
\end{proof}

The earlier sign compression, $2\lambda_2^s\le\mathbb Er^s$ and
$\lambda_1\ge1$, now gives
\[
 \gap(T_r)\ge1-[U(\tau)/2]^{1/s}>0\qquad(U(\tau)<2).
\]
There is a unique $q_0\in(0,p_*)$ with
$h(q_0)+K^s(1-q_0)=2$, since this expression decreases strictly from
infinity to a value below two. Its tail budget
$\tau_c=2-K^s(1-q_0)$ satisfies
$2-K^{s-1}<2-K^s<\tau_c<2$; $U(\tau)<2$ exactly when
$g\le\tau<\tau_c$. For $s=2$ direct expansion gives
\[
 q_0=\frac{(1-K)^2}{2-2K+K^2},\qquad
 \tau_c=2-\frac{K^2}{2-2K+K^2}.
\]
In particular $K=3/4$ gives $\tau_c=25/17$, improving the elementary
$5/4$. This is only a sufficient spectral bound: a scalar extremizer
with total moment two need not saturate the two eigenvalue inequalities.
The moment optimization is an application of classical conditional
Jensen and extremal-law reasoning; see~\cite{BertsimasPopescu} for the
general method, without an assertion of priority for this formula.

\section{The universal strict-tail threshold}\label{sec:universal-tail}
The missing link is a necessary escaping-moment estimate for an arbitrary
collapsing family. Reverse hypercontractivity forces one normalized sign
part of a second eigenvector to have a small low quasi-norm. Truncation
then forces its energy into the high-density region. This proves the
additional unit of tail cost, without assuming the constructed density
floor or an a priori decomposition into background and tilt.

\begin{theorem}\label{thm:universal-tail}
Fix the parameters above. For every $g\le\tau<1+g$ there is an explicit
constant $\gamma_0>0$, depending only on $\rho,K,\tau$, such that every
nonnegative mean-one density on every finite uniform Boolean cube with
$\mathbb E[r^s\mathbf1_{r>K}]\le\tau$ satisfies
$\gap(T_r)>\gamma_0$. No positive lower bound on $r$ is required.
For $\tau<g$ the class is empty. At $\tau=1+g$, and hence at every larger
budget, Theorem~\ref{thm:attained-cutoff} gives a positive family with
floor $K$ and gap tending to zero. Thus $1+g$ is the exact universal
strict-tail collapse threshold for $k_s<K<1$.
\end{theorem}
\begin{proof}
Write $B=\tau+K^s$, $p=1-\rho$, $q=1+\rho$ and
$\theta=1/s^2$. Classical forward and reverse Boolean
hypercontractivity~\cite[Theorem~3.1 and Corollary~3.3]{MORSS} give
\begin{align}
 \langle F,P_\rho F\rangle
   &=\|P_{\sqrt\rho}F\|_2^2\le\|F\|_q^2,\label{eq:317-forward}\\
 \langle F,P_\rho G\rangle&\ge\|F\|_p\|G\|_p
                 \quad(F,G\ge0).\label{eq:317-reverse}
\end{align}
The reverse inequality has parameter condition
$\rho\le\sqrt{(1-p)(1-p)}$, with equality here, and permits zeros.
The quantities at exponent $p<1$ are quasi-norms; no triangle inequality
for them will be used. H\"older with $q/(2-q)=s$ gives
\[
 \|\sqrt z\,h\|_q^2\le(\mathbb Ez^s)^{1/s}\|h\|_2^2
 \quad(z\ge0).
\]
Since $\mathbb Er^s\le B$, this implies $\lambda_1\le B^{1/s}$.
The unit vector $\sqrt r$ has Rayleigh quotient
$\langle r,P_\rho r\rangle\ge(\mathbb Er)^2=1$, so
$\lambda_1\ge1$.

Put $\gamma=1-\lambda_2/\lambda_1$. If $\lambda_2=0$, the gap is one
and the assertion is immediate. Otherwise work on the active set
$\{r>0\}$. The restricted matrix has strictly positive entries and a
strictly positive Perron vector. A normalized real second eigenvector
vanishes off this set, is orthogonal to the Perron vector, and therefore
has both signs. Write it as $v=f_+-f_-$, with disjoint nonnegative parts,
$a=\|f_+\|_2>0$, $b=\|f_-\|_2>0$, $a^2+b^2=1$, and
$c=\langle f_+,T_rf_-\rangle>0$. The eigenvector equation implies
\[
 \langle f_+,T_rf_+\rangle=\lambda_2a^2+c,\qquad
 \langle f_-,T_rf_-\rangle=\lambda_2b^2+c.
\]
The compression to the orthonormal pair $u=f_+/a,w=f_-/b$ is
\[
 \lambda_2I+c\begin{pmatrix}a^{-2}&(ab)^{-1}\\(ab)^{-1}&b^{-2}\end{pmatrix}.
\]
Its upper eigenvalue is $\lambda_2+c/(a^2b^2)\le\lambda_1$.
Thus the \emph{normalized} cross term obeys
\[
 d=\langle u,T_rw\rangle=c/(ab)
       \le(\lambda_1-\lambda_2)ab\le B^{1/s}\gamma/2.
\]
Both sign Rayleigh quotients are at least $\lambda_2$.
Apply~\eqref{eq:317-reverse} to $\sqrt r\,u$ and $\sqrt r\,w$.
Choose the sign vector $h$ with the smaller $p$ quasi-norm after
multiplication by $\sqrt r$. For $F=\sqrt r\,h$ we obtain
\begin{equation}\label{eq:317-sign}
 \|F\|_p\le(B^{1/s}\gamma/2)^{1/2},\qquad
 \langle F,P_\rho F\rangle\ge\lambda_2\ge1-\gamma.
\end{equation}
The argument tolerates arbitrarily unbalanced sign norms $a,b$.

Fix $M>K$, let $F_0=F\mathbf1_{r\le M}$,
$F_1=F\mathbf1_{r>M}$, and $H_M=\mathbb E[r^s\mathbf1_{r>M}]$.
One checks $1/q=\theta/p+(1-\theta)/2$. Ordinary H\"older applied to
$Z^q=(Z^p)^{\theta q/p}(Z^2)^{(1-\theta)q/2}$ proves
\[
 \|Z\|_q\le\|Z\|_p^\theta\|Z\|_2^{1-\theta}.
\]
The two powers are positive and sum to one, so this use of H\"older has
exponents greater than one even though $p<1$.
Since $\|F_0\|_2\le\sqrt M$ and $\|F_0\|_p\le\|F\|_p$,
\eqref{eq:317-forward} and~\eqref{eq:317-sign} yield
\[
 E_0:=\langle F_0,P_\rho F_0\rangle
       \le M^{1-\theta}(B^{1/s}\gamma/2)^\theta.
\]
The same forward estimate and H\"older yield
$E_1:=\langle F_1,P_\rho F_1\rangle\le H_M^{1/s}$.
The Hilbert-space triangle after applying $P_{\sqrt\rho}$ gives
$\sqrt{\langle F,P_\rho F\rangle}\le\sqrt{E_0}+\sqrt{E_1}$.
Consequently
\begin{equation}\label{eq:317-escape}
 H_M\ge\left[\sqrt{1-\gamma}
 -M^{(1-\theta)/2}(B^{1/s}\gamma/2)^{\theta/2}\right]_+^{2s}.
\end{equation}
For a collapsing family with fixed $B$, this implies
$\liminf H_M\ge1$ for each fixed $M$.

Put $A=C_sK^{s-1}$, so $g=A(1-K)$. The pointwise scalar inequality
\[
 z^s\mathbf1_{z>K}\ge A(z-K)\qquad(z\ge0)
\]
holds because its right side is nonpositive below $K$, whereas
$z^s/(z-K)$ has its global minimum $A$ at $z=sK/(s-1)$ above $K$.
Apply it to $z=r\mathbf1_{r\le M}$ on the original probability space.
Since $\mathbb E[r\mathbf1_{r>M}]\le B/M^{s-1}$, the middle strict
tail costs at least $g-AB/M^{s-1}$. It is disjoint from the high tail,
so
\begin{equation}\label{eq:317-budget}
 \tau\ge H_M+g-AB/M^{s-1}.
\end{equation}
For any collapsing family, first fix $M$, use~\eqref{eq:317-escape},
and then send $M\to\infty$ in~\eqref{eq:317-budget}. This gives
$\tau\ge1+g$. No density floor or exchange of moment limits is involved.

For an explicit positive bound, let $\varepsilon=1+g-\tau\in(0,1]$ and
\begin{align*}
 M&=\max\{K+1,(4AB/\varepsilon)^{1/(s-1)}\},\\
 D&=1-(1-3\varepsilon/4)^{1/(2s)},\\
 \gamma_0&=\min\{D/4,\,
 2B^{-1/s}(D/4)^{2/\theta}M^{-(1-\theta)/\theta}\}.
\end{align*}
These constants are finite and positive, with $\gamma_0<1$.
If $\gamma\le\gamma_0$, the subtracted core term in
\eqref{eq:317-escape} is at most $D/4$ and
$\sqrt{1-\gamma}\ge1-\gamma\ge1-D/4$. Hence
\[
 H_M\ge(1-D/2)^{2s}>(1-D)^{2s}=1-3\varepsilon/4.
\]
But~\eqref{eq:317-budget} gives
$H_M\le\tau-g+\varepsilon/4=1-3\varepsilon/4$, a contradiction.
Thus $\gamma>\gamma_0$. The empty regime $\tau<g$ was proved in
Theorem~\ref{thm:scalar-envelope}; equality $\tau=g$ is included whenever
the finite-cube class exists. Finally Theorem~\ref{thm:attained-cutoff}
uses exactly this $\rho,s,K$ and strict event, and supplies collapse at
$1+g$ with floor $K$. This completes the sharpness assertion.
\end{proof}

For $s=2$, $\rho=1/3$ and $K=3/4$, the exact threshold is
$1+4(1/4)(3/4)=7/4$. The intermediate scalar bound $25/17$ remains
valid and can be numerically much simpler than the conservative explicit
constant above. Neither scalar optimality nor the construction alone
proved the universal threshold: the necessary escape
estimate~\eqref{eq:317-escape} closes that gap. External priority of this
specific synthesis is unestablished. Its inputs are classical; related
transmission of log-Sobolev and spectral-gap inequalities for perturbed
diffusion measures is studied, for example, in~\cite{Cattiaux}.
Those generator/measure results are not silently identified with the
finite transfer-matrix statement here. This theorem does not add marginal
matching, finite-range interactions or a physical interpretation.
